\subsection{Grouped Repeat Prefix REPG}
REPG repeats an ENDG-terminated block of eligible instructions with a D-register
counter. The first grouped instruction carries the REPG prefix byte; the final
grouped instruction carries the ENDG prefix byte. Instructions between those
endpoints are ordinary instruction records with their own word 0 length fields.
The grouped instructions execute in program order and have no condition-code
selector.

Architecturally, REPG is a compact spelling for a counted scalar loop over a
short straight-line instruction trace. One group iteration executes every
instruction from the REPG-marked instruction through the ENDG-marked instruction
exactly once, in instruction-stream order. After a successful group iteration,
the selected D register is moved one step toward zero. A zero counter performs
no group iteration. Positive counters count down; negative counters count up
toward zero.

\paragraph{Group Discovery.}
A REPG group is discovered entirely from the architectural instruction stream.
When a REPG prefix is seen, instruction interpretation follows the explicit word
0 length fields of subsequent instruction records in the same naturally aligned
64-byte grouping window until it reaches an instruction carrying ENDG. ENDG may
appear on the same instruction as REPG, which forms a one-instruction group. A
well-formed REPG group has exactly one active REPG start marker, exactly one ENDG
terminator, and no nested repeat prefix inside the group.

\paragraph{Fast-Contract Source Form.}
\texttt{REPGF} is a source-level assertion that uses the same architectural
encoding as \texttt{REPG}. It does not allocate a second prefix byte and does
not force a particular implementation strategy. It marks the group as a
REPG-fast candidate: the selected counter register is outside the group writes,
the effective addresses are regular, and the group avoids control, state, and
nested-repeat operations. Runtime properties such as ordinary mapped memory
remain architectural preconditions of the fast class rather than extra encoded
bits.

\paragraph{Grouping Window.}
The complete byte range occupied by the grouped instructions must fit within the
naturally aligned 64-byte grouping window that contains the first grouped
instruction. The first instruction does not need to be aligned to the beginning
of the window; only containment matters. This keeps architectural group
recognition bounded while avoiding any dependency on a particular
implementation's instruction-fetch granularity.

\paragraph{Execution and Faults.}
Each grouped instruction retains its ordinary operand decoding,
effective-address calculation, privilege checks, FLAGS or FFLAGS behavior, and
memory side effects. Address-update operands update as they would in scalar
execution for each completed instruction.

If a fault occurs, completed iterations and completed instructions in the
current iteration remain committed, and the fault is precise at the faulting
grouped instruction. The counter update for an iteration occurs only after the
whole group iteration completes successfully. A REPG fault frame saves the PC of
the faulting grouped instruction and a compact repeat continuation in
\texttt{FAULT\_AUX}. The group start is not stored as a full address; it is
reconstructed from a five-bit unsigned negative word displacement:

\[
  \texttt{group\_start\_pc}
  =
  \texttt{fault\_pc}
  -
  2 \times \texttt{group\_start\_delta\_words}
\]

The group index is not saved. If the handler returns without editing the frame,
IRET resumes at the saved PC under the active REPG continuation and then
proceeds through the rest of the group.

\paragraph{Legality.}
A REPG group is intended to be straight-line code. Branches, calls, returns,
traps, system-call returns, interrupt returns, atomics, stack operations,
nested repeat prefixes, and state-mutating system, TLB, cache, or fence
operations are outside the REPG group model.

REPG optimization is divided into REPG-fast and REPG-general. REPG-fast
identifies groups that are suitable for internal SIMD or streaming recognition.
REPG-general preserves exact scalar semantics for less regular groups, including
groups with selected-counter dependencies.

\begingroup\small
\begin{longtable}{@{}p{1.35in}p{4.45in}@{}}
\toprule
\textbf{Item} & \textbf{Value}\\
\midrule
Syntax & source uses \texttt{REPG Dn,} followed by a braced instruction list\\
Fast-contract syntax & source may use \texttt{REPGF Dn,} followed by a braced instruction list; the emitted encoding is the same as \texttt{REPG}\\
Condition selector & none\\
Counter & D register\\
Group termination & ENDG prefix on the final grouped instruction\\
Eligible instructions & repeatable operations allowed by the instruction metadata\\
Streaming candidate marker & REPG-fast candidate when the group satisfies the fast rules\\
Grouping window & grouped instruction bytes must fit in one naturally aligned 64-byte window\\
Commit rule & completed iterations and completed prior group instructions commit\\
Fault continuation & saved PC is the faulting grouped instruction; group-start delta is a 5-bit negative word displacement; no group index is saved\\
FFLAGS & OR of completed iterations\\
\bottomrule
\end{longtable}
\manualtablecaption{REPG Prefix Rules}
\endgroup

\begingroup\small
\begin{longtable}{@{}p{0.95in}p{2.15in}p{2.55in}@{}}
\toprule
\textbf{Class} & \textbf{Availability} & \textbf{Implementation}\\
\midrule
\texttt{REPG-fast} &
\begin{tabular}[t]{@{}l@{}}
counter Dn is not written\\
regular effective addresses\\
ordinary memory\\
no control or state changes
\end{tabular}
& internal SIMD or streaming optimization is possible\\
\texttt{REPG-general} &
\begin{tabular}[t]{@{}l@{}}
counter Dn may be read or written\\
state-query instructions may appear\\
wait-like instructions may appear\\
examples: WAIT, CPUID
\end{tabular}
& exact scalar semantics, with no fast streaming guarantee\\
\bottomrule
\end{longtable}
\manualtablecaption{REPG Optimization Classes}
\endgroup

\paragraph{Example.}
\begin{quote}
\begingroup\small\ttfamily
\begin{tabular}{@{}l@{}}
REPG D2, \{\\
\quad MOV.L  [A1++], D4\\
\quad MADD.L [A0++], D4, D1\\
\}\\
\end{tabular}
\endgroup
\end{quote}

The example repeats a two-instruction dot-product step. Each iteration loads one
source value, performs one multiply-add using the other source stream, advances
the update-eligible addresses, and then updates D2 toward zero after the whole
group completes.
